% GENERATED FILE. Edit canonical lessons, then run npm run generate:books.
% canonical-source-hash: fnv1a64:22c9152342931a7e
% canonical-lessons: MR-C448-sandpani, MR-C448-durgandhi, MR-C448-pradusit, MR-C448-kujne, MR-C448-sadne

\chapter{Sewage and Stench}
\label{ch:mr-sewage-and-stench}
\hlchaptermodality{448}

\begin{chapteropening}
\textbf{By the end of this chapter:} \emph{I can talk about sewage, a stench and polluted water, and say things rot and decay.}

Complete the last lesson of chapter 448: I can talk about sewage, a stench and polluted water, and say things rot and decay.
\end{chapteropening}

\section[\texorpdfstring{\mr{सांडपाणी} (sāṇḍpāṇī)}{सांडपाणी (sāṇḍpāṇī)}]{\mr{सांडपाणी} (sāṇḍpāṇī) — waste water; sewage}

\label{lesson:MR-C448-sandpani}



\begin{quote}
Before the new one: say the Marathi for to chirp, then the Marathi for migration.
\end{quote}

\subsection*{You'll want to know: \mr{सांडपाणी}}
\textbf{\mr{सांडपाणी}} — \emph{sāṇḍpāṇī} — \textquotedblleft{}waste water; sewage\textquotedblright{}.

\begin{quote}
\mr{सांडपाणी} \mr{नदीत} \mr{सोडू} \mr{नये}.
\end{quote}

\begin{quote}
\emph{sāṇḍpāṇī nadīt soḍū naye.}
\end{quote}

\begin{quote}
Waste water should not be let into the river.
\end{quote}

Neuter.

Used water from homes and factories, which should be cleaned before it reaches a river.

\begin{cousinweb}
\textbf{\mr{सांडपाणी}} (\emph{sāṇḍpāṇī}) is built from \emph{sāṇḍ-}, \textquotedblleft{}spilt, let go\textquotedblright{}, and \textbf{\mr{पाणी}}, \textquotedblleft{}water\textquotedblright{}, inherited from Sanskrit \emph{pānīya} (CDIAL 8082); the first part is not settled.
\end{cousinweb}

\begin{grammarlens}[title={What you've built}]
One more word for getting around.
\end{grammarlens}

\subsection*{Your turn}
\begin{itemize}
\raggedright
  \item \emph{Say it:} \emph{sāṇḍpāṇī}
  \item \emph{Say it:} \emph{sāṇḍpāṇī}, once more
  \item \emph{Say it:} \emph{sāṇḍpāṇī nadīt soḍū naye.} — Waste water should not be let into the river.
  \item \emph{Recall:} say the Marathi for to chirp, then the Marathi for migration, then say \emph{sāṇḍpāṇī} again
\end{itemize}

\begin{tcolorbox}[breakable,colback=teal!4,colframe=teal!35!black,title={Before you move on}]
What does \mr{सांडपाणी} mean? (\textquotedblleft{}Waste water; sewage\textquotedblright{}.) Say it once more.
\end{tcolorbox}
\section[\texorpdfstring{\mr{दुर्गंधी} (durgandhī)}{दुर्गंधी (durgandhī)}]{\mr{दुर्गंधी} (durgandhī) — a stench}

\label{lesson:MR-C448-durgandhi}



\begin{quote}
Before the new one: say the Marathi for migration, then the Marathi for waste water.
\end{quote}

\subsection*{You'll want to know: \mr{दुर्गंधी}}
\textbf{\mr{दुर्गंधी}} — \emph{durgandhī} — \textquotedblleft{}a stench\textquotedblright{}.

\begin{quote}
\mr{सांडपाण्यामुळे} \mr{गल्लीत} \mr{दुर्गंधी} \mr{पसरली}.
\end{quote}

\begin{quote}
\emph{sāṇḍpāṇyāmuḷe gallīt durgandhī pasarlī.}
\end{quote}

\begin{quote}
A stench spread through the lane because of the sewage.
\end{quote}

Feminine.

A feminine noun; the opposite is \mr{सुगंध}, a fragrance.

\begin{cousinweb}
\textbf{\mr{दुर्गंधी}} (\emph{durgandhī}) is a tatsama, taken from Sanskrit \emph{durgandha}, \textquotedblleft{}a bad smell, a stink\textquotedblright{} (Monier-Williams).
\end{cousinweb}

\begin{grammarlens}[title={What you've built}]
One more word for getting around.
\end{grammarlens}

\subsection*{Your turn}
\begin{itemize}
\raggedright
  \item \emph{Say it:} \emph{durgandhī}
  \item \emph{Say it:} \emph{durgandhī}, once more
  \item \emph{Say it:} \emph{sāṇḍpāṇyāmuḷe gallīt durgandhī pasarlī.} — A stench spread through the lane because of the sewage.
  \item \emph{Recall:} say the Marathi for migration, then the Marathi for waste water, then say \emph{durgandhī} again
\end{itemize}

\begin{tcolorbox}[breakable,colback=teal!4,colframe=teal!35!black,title={Before you move on}]
What does \mr{दुर्गंधी} mean? (\textquotedblleft{}A stench\textquotedblright{}.) Say it once more.
\end{tcolorbox}
\section[\texorpdfstring{\mr{प्रदूषित} (pradūṣit)}{प्रदूषित (pradūṣit)}]{\mr{प्रदूषित} (pradūṣit) — polluted}

\label{lesson:MR-C448-pradusit}



\begin{quote}
Before the new one: say the Marathi for waste water, then the Marathi for a stench.
\end{quote}

\subsection*{You'll want to know: \mr{प्रदूषित}}
\textbf{\mr{प्रदूषित}} — \emph{pradūṣit} — \textquotedblleft{}polluted\textquotedblright{}.

\begin{quote}
\mr{शहरातली} \mr{हवा} \mr{प्रदूषित} \mr{झाली} \mr{आहे}.
\end{quote}

\begin{quote}
\emph{śahrātlī havā pradūṣit jhālī āhe.}
\end{quote}

\begin{quote}
The air in the city has become polluted.
\end{quote}

It does not change: \mr{प्रदूषित} \mr{हवा}, \mr{प्रदूषित} \mr{पाणी}.

\begin{cousinweb}
\textbf{\mr{प्रदूषित}} (\emph{pradūṣit}) is a tatsama, taken from Sanskrit \emph{pradūṣita}, \textquotedblleft{}corrupted, spoilt\textquotedblright{} (Monier-Williams).
\end{cousinweb}

\begin{grammarlens}[title={What you've built}]
One more word for getting around.
\end{grammarlens}

\subsection*{Your turn}
\begin{itemize}
\raggedright
  \item \emph{Say it:} \emph{pradūṣit}
  \item \emph{Say it:} \emph{pradūṣit}, once more
  \item \emph{Say it:} \emph{śahrātlī havā pradūṣit jhālī āhe.} — The air in the city has become polluted.
  \item \emph{Recall:} say the Marathi for waste water, then the Marathi for a stench, then say \emph{pradūṣit} again
\end{itemize}

\begin{tcolorbox}[breakable,colback=teal!4,colframe=teal!35!black,title={Before you move on}]
What does \mr{प्रदूषित} mean? (\textquotedblleft{}Polluted\textquotedblright{}.) Say it once more.
\end{tcolorbox}
\section[\texorpdfstring{\mr{कुजणे} (kujṇe)}{कुजणे (kujṇe)}]{\mr{कुजणे} (kujṇe) — to rot}

\label{lesson:MR-C448-kujne}



\begin{quote}
Before the new one: say the Marathi for a stench, then the Marathi for polluted.
\end{quote}

\subsection*{You'll want to know: \mr{कुजणे}}
\textbf{\mr{कुजणे}} — \emph{kujṇe} — \textquotedblleft{}to rot\textquotedblright{}.

\begin{quote}
\mr{पावसात} \mr{लाकूड} \mr{कुजलं}.
\end{quote}

\begin{quote}
\emph{pāvsāt lākūḍ kujlaṁ.}
\end{quote}

\begin{quote}
The wood rotted in the rain.
\end{quote}

Past (he …): \mr{तो} \mr{कुजला}, \emph{to kujlā}.

Wood, leaves and fruit \mr{कुजतात} slowly.

\begin{cousinweb}
\textbf{\mr{कुजणे}} (\emph{kujṇe}) is inherited: CDIAL lists Marathi \emph{kuj̈ṇẽ}, \textquotedblleft{}to rot\textquotedblright{}, under \emph{*kujjati}, \textquotedblleft{}rots\textquotedblright{} (CDIAL 3222).
\end{cousinweb}

\begin{grammarlens}[title={What you've built}]
One more word for getting around.
\end{grammarlens}

\subsection*{Your turn}
\begin{itemize}
\raggedright
  \item \emph{Say it:} \emph{kujṇe}
  \item \emph{Say it:} \emph{kujṇe}, once more
  \item \emph{Say it:} \emph{pāvsāt lākūḍ kujlaṁ.} — The wood rotted in the rain.
  \item \emph{Recall:} say the Marathi for a stench, then the Marathi for polluted, then say \emph{kujṇe} again
\end{itemize}

\begin{tcolorbox}[breakable,colback=teal!4,colframe=teal!35!black,title={Before you move on}]
What does \mr{कुजणे} mean? (\textquotedblleft{}To rot\textquotedblright{}.) Say it once more.
\end{tcolorbox}
\section[\texorpdfstring{\mr{सडणे} (saḍṇe)}{सडणे (saḍṇe)}]{\mr{सडणे} (saḍṇe) — to decay}

\label{lesson:MR-C448-sadne}



\begin{quote}
Before the new one: say the Marathi for polluted, then the Marathi for to rot.
\end{quote}

\subsection*{You'll want to know: \mr{सडणे}}
\textbf{\mr{सडणे}} — \emph{saḍṇe} — \textquotedblleft{}to decay\textquotedblright{}.

\begin{quote}
\mr{टोपलीतले} \mr{आंबे} \mr{सडले}.
\end{quote}

\begin{quote}
\emph{ṭoplītle āmbe saḍle.}
\end{quote}

\begin{quote}
The mangoes in the basket went rotten.
\end{quote}

Past (he …): \mr{तो} \mr{सडला}, \emph{to saḍlā}.

Sharper than \mr{कुजणे}: it often smells.

\begin{cousinweb}
\textbf{\mr{सडणे}} (\emph{saḍṇe}) is inherited: CDIAL lists Marathi \emph{saḍṇẽ}, \textquotedblleft{}to rot\textquotedblright{}, under Sanskrit \emph{śaṭati} (CDIAL 12268).
\end{cousinweb}

\begin{grammarlens}[title={What you've built}]
One more word for getting around.
\end{grammarlens}

\subsection*{Your turn}
\begin{itemize}
\raggedright
  \item \emph{Say it:} \emph{saḍṇe}
  \item \emph{Say it:} \emph{saḍṇe}, once more
  \item \emph{Say it:} \emph{ṭoplītle āmbe saḍle.} — The mangoes in the basket went rotten.
  \item \emph{Recall:} say the Marathi for polluted, then the Marathi for to rot, then say \emph{saḍṇe} again
\end{itemize}

\begin{tcolorbox}[breakable,colback=teal!4,colframe=teal!35!black,title={Before you move on}]
What does \mr{सडणे} mean? (\textquotedblleft{}To decay\textquotedblright{}.) Say it once more.
\end{tcolorbox}
